% !PRINTED 91
\weberpage{91}

%<<CONT
we can also in these cases take $e_{1,h}$ so that $c_{1,h}\,e_{1,h}$
does not become divisible by $q_h$.

The requirements herewith imposed on the $e_{k,h}$ thus consist only in
this, that they are to be divisible by certain primes and not divisible
by others, and so they are still compatible with one another in
infinitely many ways. But this last assumption about the $e_{k,h}$ had
only the purpose of showing that the earlier, more general requirements
can always be satisfied. There may still be many other ways of meeting
them.

\noindent\rule{4cm}{0.4pt}

\begin{center}
\gesperrt{Survey of the Results.}
\end{center}

If it is a question of representing all cyclotomic fields, and each one
only once, and indeed in the simplest possible way, by cyclotomic
periods, one proceeds as follows:

Take an arbitrary system of prime powers as invariants

\[
i_1, i_2, \dots i_v
\]

and then choose a number $m$ according to the following conditions.

If $p^\pi$ is the highest power of a prime $p$ that occurs among the
invariants, then take $p$ into $m$ at most $(\pi+1)$ times, or, if
$p = 2$, $(\pi+2)$ times.

A prime $q$ that does not divide the invariants at all may be taken
into $m$ only singly; but $m$ may have only such single prime factors
$q$ for which $q - 1$ is divisible by one of the prime factors $p$ of
the invariants. Hence $p$ itself, too, can occur singly in $m$ only if
$p - 1$ is divisible by one of the other $p$.

Further, $m$ must also satisfy the condition that among the index
moduli

\[
c_1, c_2, \dots c_\mu
\]

$\varrho$ distinct ones, say $c_1, c_2, \dots c_\varrho$, are divisible
by $i_1, i_2, \dots i_\varrho$, whenever $i_1, i_2, \dots i_\varrho$
are powers of one and the same prime. Then determine the quantities
$i_{k,h}$ and $c_{k,h}$ by the formulae (11) and divide the given
invariants into systems, so that all powers of one and the same prime
are united in one system.

For each of these systems form the matrix (24) and choose the numbers
$c_{k,h}$ so that from every such matrix a
% [sic: Weber prints c_{k,h}; these quantities were fixed by (11), whereas the matrix entries and the ensuing construction require choosing e_{k,h}.]
%CONT>>
